\documentclass[a4paper,11pt,oneside,openany]{book}
\usepackage{formula_style}

\title{CSR102: Design Thinking and Complex Problem Solving\\Frameworks, Canvases \& Quick Reference Handbook}
\author{Academic Review Board \& Curriculum Development Cell}
\date{University Academic Edition 2026}

\begin{document}

\frontmatter

\begin{titlepage}
    \centering
    \vspace*{2cm}
    {\Huge\bfseries\color{LPUNavy} CSR102: DESIGN THINKING AND COMPLEX PROBLEM SOLVING\par}
    \vspace{0.8cm}
    {\Large\bfseries\color{LPUOrange} Frameworks, Canvases, Methodologies \& Quick Reference Handbook\par}
    \vspace{0.5cm}
    {\large The Complete Visual Reference Companion for Engineering Design Practice\par}
    \vspace{2cm}
    \begin{tcolorbox}[colback=FormulaBg,colframe=LPUNavy,arc=3mm,boxrule=1pt,width=0.88\textwidth]
        \centering\bfseries\color{LPUNavy}
        CORE HANDBOOK SECTIONS \& TOOLKITS:\\[2mm]
        \normalfont\small
        $\bullet$ \textbf{Unit 1:} Stanford d.school Model, Double Diamond, Wicked Problems Matrix, Cynefin Framework\\[1mm]
        $\bullet$ \textbf{Unit 2:} Stakeholder Power-Interest Grid, Empathy Map, Customer Journey Map, 5-Whys, Fishbone, POV/HMW\\[1mm]
        $\bullet$ \textbf{Unit 3:} Reasoning Typology, 12 Cognitive Biases \& Debiasing, MCDA Scoring Formulas, MECE Decision Trees\\[1mm]
        $\bullet$ \textbf{Unit 4:} SCAMPER Matrix, Brainwriting 6-3-5, DFV Venn Scoring, Pugh Decision Matrix, 2$\times$2 Prioritization\\[1mm]
        $\bullet$ \textbf{Unit 5:} Service Blueprint Architecture, MVP Experiment Typology, Usability Metrics (SUS, TSR, SEQ), Feedback Grid\\[1mm]
        $\bullet$ \textbf{Unit 6:} Barbara Minto SCQA Pyramid, Executive Narrative Pitch Arc, Kolb \& Gibbs Reflective Cycles
    \end{tcolorbox}
    \vfill
    {\large Lovely Professional University Open Academic Resources\par}
    {\small Course Code: CSR102 \textbullet\ School of Computer Science and Engineering\par}
    \vspace*{1.5cm}
\end{titlepage}

\tableofcontents
\newpage

\mainmatter

\chapter{Unit 1: Foundations of Problem Solving and Design Thinking}
\label{chap:unit1-frameworks}

\section{Problem Typology: Well-Structured vs. Ill-Structured vs. Wicked}

\begin{frameworkbox}{The Three Classes of Engineering \& Socio-Technical Problems}
\begin{tabularx}{\linewidth}{|l|X|X|X|} \hline
\textbf{Dimension} & \textbf{Well-Structured (Tame)} & \textbf{Ill-Structured (Complex)} & \textbf{Wicked Problem} \\ \hline
\textbf{Problem Definition} & Clear, complete, unambiguous. Known initial and goal states. & Vaguely defined; conflicting descriptions; missing variables. & No definitive formulation; framing the problem is the problem. \\ \hline
\textbf{Stopping Rule} & Explicit algorithm terminates when criterion is satisfied. & Pragmatic termination governed by resources, time, budget. & No stopping rule; process ends when resources or patience run out. \\ \hline
\textbf{Solution Testing} & True/False; Correct/Incorrect testable against proofs. & Good/Bad; Better/Worse evaluated across multiple criteria. & No immediate or ultimate test; solutions ripple uncontrollably. \\ \hline
\textbf{Reversibility} & Trial-and-error possible in simulation or sandbox. & Costly to reverse; operational inertia. & Every intervention is a ``one-shot operation''; irreversible consequences. \\ \hline
\textbf{Archetype Example} & Sorting algorithm, solving linear equations, bridge load stress. & Redesigning hospital queue management, reducing bug latency. & Climate change, healthcare inequality, urban homelessness. \\ \hline
\end{tabularx}
\end{frameworkbox}

\begin{memorytip}
\textbf{Rittel \& Webber (1973) 10 Hallmarks of Wicked Problems:}\\
1. No definitive formulation. 2. No stopping rule. 3. Solutions are not true-or-false, but good-or-bad. 4. No immediate/ultimate test. 5. Every solution is a one-shot operation. 6. No enumerable set of solutions. 7. Every wicked problem is essentially unique. 8. Every wicked problem is a symptom of another problem. 9. Discrepancies can be explained in numerous ways. 10. The planner has no right to be wrong (high stakes).
\end{memorytip}

\section{The Stanford d.school 5-Stage Iterative Design Framework}

\begin{frameworkbox}{The 5-Stage Non-Linear Process Flow}
\begin{enumerate}[label=\textbf{\arabic*.}, leftmargin=2em]
    \item \textbf{Empathize (Observe, Engage, Immerse):}
    \begin{itemize}
        \item \textit{Objective:} Understand user experiences, motivations, emotional triggers, and unarticulated unmet needs.
        \item \textit{Key Methods:} Semi-structured interviews, shadowing/contextual inquiry, extreme user studies, fly-on-the-wall observation.
    \end{itemize}
    \item \textbf{Define (Unpack, Synthesize, Frame):}
    \begin{itemize}
        \item \textit{Objective:} Process empathy data into a human-centered, actionable Point-of-View (POV) problem statement.
        \item \textit{Key Methods:} Affinity mapping, empathy maps, persona archetypes, user journey mapping, HMW formulation.
    \end{itemize}
    \item \textbf{Ideate (Diverge, Explore, Expand):}
    \begin{itemize}
        \item \textit{Objective:} Generate vast quantities and diverse varieties of solution possibilities before premature selection.
        \item \textit{Key Methods:} Brainwriting 6-3-5, SCAMPER, Crazy Eights, analogical thinking, worst possible idea.
    \end{itemize}
    \item \textbf{Prototype (Build to Think, Make Concrete):}
    \begin{itemize}
        \item \textit{Objective:} Transform concepts into tangible artifacts that provoke questions, test assumptions, and expose flaws cheaply.
        \item \textit{Key Methods:} Paper prototypes, cardboard mockups, digital wireframes, Wizard of Oz, service walkthroughs.
    \end{itemize}
    \item \textbf{Test (Validate, Learn, Iterate):}
    \begin{itemize}
        \item \textit{Objective:} Observe real users interacting with the prototype in authentic contexts to refine solutions and reframe problems.
        \item \textit{Key Methods:} Think-aloud usability testing, task completion tracking, feedback capture grids, A/B split testing.
    \end{itemize}
\end{enumerate}
\end{frameworkbox}

\begin{memorytip}
\textbf{Critical Concept:} Design Thinking is \textbf{non-linear}. Testing frequently leads back to Empathize (discovering user insights) or Define (reframing the core problem). It is an iterative learning engine, not a sequential waterfall phase-gate.
\end{memorytip}

\section{The UK Design Council Double Diamond Model}

\begin{frameworkbox}{Divergent \& Convergent Thinking Cycles}
\begin{itemize}
    \item \textbf{Diamond 1: Finding the Right Problem (Problem Space)}
    \begin{itemize}
        \item \textit{Discover (Divergence):} Broaden perspective, gather broad qualitative and contextual data, question initial assumptions.
        \item \textit{Define (Convergence):} Synthesize findings, filter noise, focus on the most impactful core root challenge.
    \end{itemize}
    \item \textbf{Diamond 2: Developing the Right Solution (Solution Space)}
    \begin{itemize}
        \item \textit{Develop (Divergence):} Brainstorm multifaceted solution concepts, cross-pollinate ideas across domains.
        \item \textit{Deliver (Convergence):} Down-select, prototype, test, eliminate weak solutions, and implement the validated output.
    \end{itemize}
\end{itemize}
\end{frameworkbox}

\section{The Cynefin Framework for Decision Contexts}

\begin{frameworkbox}{Dave Snowden's Sense-Making Framework}
\begin{tabularx}{\linewidth}{|l|l|l|X|} \hline
\textbf{Domain} & \textbf{System Characteristic} & \textbf{Response Cycle} & \textbf{Engineering Action} \\ \hline
\textbf{Clear (Simple)} & Predictable, cause-and-effect direct. & Sense $\to$ Categorize $\to$ Respond & Apply Best Practice, checklists, standard operating procedures. \\ \hline
\textbf{Complicated} & Multiple valid answers, cause-and-effect separated. & Sense $\to$ Analyze $\to$ Respond & Apply Good Practice, subject matter expert analysis, trade-off matrices. \\ \hline
\textbf{Complex} & Emergent, unpredictably adaptive, no repeated patterns. & Probe $\to$ Sense $\to$ Respond & Deploy Safe-to-Fail experiments, rapid prototypes, iterative adaptation. \\ \hline
\textbf{Chaotic} & High turbulence, immediate crisis, no cause-and-effect link. & Act $\to$ Sense $\to$ Respond & Establish immediate stabilization, decisive triage, shift to complex. \\ \hline
\textbf{Aporetic / Confused} & Unclear which domain applies. & Deconstruct $\to$ Partition & Break system into sub-components and assign to domains. \\ \hline
\end{tabularx}
\end{frameworkbox}

\chapter{Unit 2: Problem Definition and Analysis Toolkits}
\label{chap:unit2-canvases}

\section{Stakeholder Mapping \& Power-Interest Matrix}

\begin{canvasbox}{Mendelow's Power-Interest Matrix}
\begin{tabularx}{\linewidth}{|l|X|X|} \hline
 & \textbf{Low Interest} & \textbf{High Interest} \\ \hline
\textbf{High Power} & \textbf{Keep Satisfied:} Minimal engagement required for operations, but crucial veto power. Brief periodically with executive summaries. & \textbf{Manage Closely (Key Partners):} Essential co-creators. Involve actively in co-design, steering committees, and decision gates. \\ \hline
\textbf{Low Power} & \textbf{Monitor (Minimum Effort):} Track for shifts in power or interest; update via general passive channels. & \textbf{Keep Informed:} High passion, low direct leverage. Cultivate as grassroots advocates, beta testers, and qualitative feedback sources. \\ \hline
\end{tabularx}
\end{canvasbox}

\section{The Empathy Map Canvas}

\begin{canvasbox}{Dave Gray's Empathy Map Architecture}
\begin{tabularx}{\linewidth}{|X|X|} \hline
\textbf{Says (Verbatim Expressions):} & \textbf{Thinks (Internal Mental Model):} \\
$\bullet$ Exact quotes from user interviews. & $\bullet$ Beliefs, assumptions, mental models. \\
$\bullet$ Stated goals and self-described justifications. & $\bullet$ Unvoiced worries, private expectations. \\ \hline
\textbf{Does (Observable Behaviors):} & \textbf{Feels (Emotional States):} \\
$\bullet$ Physical actions, workarounds, body language. & $\bullet$ Emotions: anxiety, frustration, delight, apathy. \\
$\bullet$ Specific pathways chosen in existing workflows. & $\bullet$ Subtle non-verbal cues and physiological signs. \\ \hline
\multicolumn{2}{|l|}{\textbf{Pains (Fears, Frustrations, Obstacles, Risks):}} \\
\multicolumn{2}{|X|}{What stands between the user and success? Where does friction occur? What are the high-consequence failure modes?} \\ \hline
\multicolumn{2}{|l|}{\textbf{Gains (Hopes, Aspirations, Desired Outcomes, Measures of Success):}} \\
\multicolumn{2}{|X|}{What does success look like for the user? What would genuinely delight them beyond minimal expectations?} \\ \hline
\end{tabularx}
\end{canvasbox}

\section{Customer Journey Mapping Framework}

\begin{canvasbox}{Core Columns of a High-Fidelity Journey Map}
\begin{tabularx}{\linewidth}{|l|X|} \hline
\textbf{Journey Stage} & Chronological phases: Awareness $\to$ Research $\to$ Onboarding $\to$ Core Usage $\to$ Support $\to$ Offboarding. \\ \hline
\textbf{Touchpoints} & Channels of interaction (Mobile app, desktop kiosk, customer desk, SMS, hardware terminal). \\ \hline
\textbf{User Actions} & What the user concretely does at each discrete milestone. \\ \hline
\textbf{Thoughts \& Mental State} & Questions asked, doubts raised, expectations formed during the transition. \\ \hline
\textbf{Emotional Curve} & Bipolar emotional valence mapped visually from $+2$ (Delight) to $-2$ (Crisis/Frustration). \\ \hline
\textbf{Pain Points / Friction} & Bottlenecks, cognitive overload, latency, errors, unexpected prerequisites. \\ \hline
\textbf{Design Opportunities} & Specific HMW intervention points derived directly from emotional troughs. \\ \hline
\end{tabularx}
\end{canvasbox}

\section{Root Cause Diagnostic Toolkits}

\begin{frameworkbox}{1. The 5-Whys Method (Sakichi Toyoda)}
\textit{Iterative interrogative technique to peel away superficial symptoms and expose the fundamental process/policy breakdown.}\\
Rule: Every answer must link logically to the preceding question through demonstrable causality, not conjecture.
\end{frameworkbox}

\begin{frameworkbox}{2. Ishikawa (Fishbone / Cause-and-Effect) 6M \& 4P Frameworks}
\begin{itemize}
    \item \textbf{Manufacturing / Technical 6M:} Manpower, Methods, Machines, Materials, Measurements, Mother Nature (Environment).
    \item \textbf{Service / Knowledge Work 4P:} People, Policies, Procedures, Plant / Technology.
\end{itemize}
\end{frameworkbox}

\section{Point-of-View (POV) Formula \& ``How Might We'' (HMW) Generators}

\begin{metricbox}{The Stanford d.school POV Formula}
\[
\text{\textbf{POV Statement}} = \underbrace{[\text{Specific User / Persona}]}_{\text{Who}} + \text{ needs } \underbrace{[\text{Deep Human Need}]}_{\text{What (Verb-driven)}} + \text{ because } \underbrace{[\text{Surprising Insight}]}_{\text{Root Emotional Driver}}
\]
\end{metricbox}

\begin{frameworkbox}{HMW Expansion Strategies (From 1 POV Statement)}
\begin{enumerate}[label=\textbf{\arabic*.}, leftmargin=2em]
    \item \textbf{Amp up the good:} How might we amplify the user's existing delight during step $X$?
    \item \textbf{Remove the bad:} How might we completely eliminate the need to wait in line?
    \item \textbf{Explore the opposite:} How might we make waiting in line the most enjoyable part of the day?
    \item \textbf{Question assumptions:} How might we eliminate the physical office entirely?
    \item \textbf{Leverage an unexpected resource:} How might we use idle bystander smartphones to accelerate verification?
    \item \textbf{Analogy to another domain:} How might we make hospital check-in feel like Disney theme-park boarding?
\end{enumerate}
\end{frameworkbox}

\section{SMART Objective Setting Framework}

\begin{canvasbox}{SMART Verification Architecture}
\begin{tabularx}{\linewidth}{|l|X|l|} \hline
\textbf{Criterion} & \textbf{Evaluation Standard} & \textbf{Example Non-Compliant vs. Compliant} \\ \hline
\textbf{S - Specific} & Clear, focused, answers Who, What, Where. & ``Improve speed'' $\to$ ``Reduce patient registration time'' \\ \hline
\textbf{M - Measurable} & Quantifiable metrics and verification telemetry. & ``Make better'' $\to$ ``Reduce median time from 18 to 4 minutes'' \\ \hline
\textbf{A - Achievable} & Realistic within technical and resource limits. & ``Zero errors forever'' $\to$ ``$\le 0.5\%$ error rate'' \\ \hline
\textbf{R - Relevant} & Directly aligned with core human and business goals. & Tracking meaningless page clicks vs. successful consultations. \\ \hline
\textbf{T - Time-bound} & Unambiguous milestone and final deadlines. & ``Soon'' $\to$ ``Within 90 calendar days of pilot launch'' \\ \hline
\end{tabularx}
\end{canvasbox}

\chapter{Unit 3: Analytical and Critical Thinking Toolkits}
\label{chap:unit3-analytical-toolkit}

\section{Epistemological Reasoning Typology}

\begin{frameworkbox}{Deductive vs. Inductive vs. Abductive Reasoning in Engineering}
\begin{tabularx}{\linewidth}{|l|X|X|X|} \hline
\textbf{Dimension} & \textbf{Deductive Reasoning} & \textbf{Inductive Reasoning} & \textbf{Abductive Reasoning} \\ \hline
\textbf{Direction} & General Rule $\to$ Specific Case & Specific Observations $\to$ General Pattern & Surprising Fact $\to$ Best Plausible Hypothesis \\ \hline
\textbf{Logic} & If premises are true, conclusion \textit{must} be true (Certainty). & If premises are true, conclusion is \textit{probable} (Probability). & Creative leap to explain anomalous observation (Plausibility). \\ \hline
\textbf{Role in Design} & Verifying constraints, safety bounds, compliance rules. & Analyzing user interview trends, telemetry logs. & Generating novel solution hypotheses during ideation. \\ \hline
\textbf{Formula} & All $A$ are $B$; $x$ is $A$; $\therefore x$ is $B$. & $x_1, x_2, \dots, x_n$ are $B$; $\therefore$ All $X$ are $B$. & Fact $C$ is observed; If $H$ were true, $C$ is natural; $\therefore H$ is likely. \\ \hline
\end{tabularx}
\end{frameworkbox}

\section{The 12 Core Cognitive Biases in Complex Problem Solving}

\begin{frameworkbox}{Cognitive Biases \& Operational Debiasing Mechanisms}
\begin{tabularx}{\linewidth}{|l|X|X|} \hline
\textbf{Cognitive Bias} & \textbf{Manifestation in Problem Solving} & \textbf{Debiasing Heuristic} \\ \hline
\textbf{Confirmation Bias} & Seeking data supporting pre-existing hypothesis; discarding contrary signals. & Appoint a designated ``Red Team'' / Devil's Advocate to argue falsification. \\ \hline
\textbf{Availability Heuristic} & Overweighting recent vivid events or memorable customer anecdotes. & Demand empirical telemetry and stratified statistical sampling. \\ \hline
\textbf{Anchoring Effect} & Fixating on the first number or specification proposed in discussions. & Establish objective reference-class forecasting before negotiations. \\ \hline
\textbf{Framing Effect} & Changing decisions based solely on whether choices are framed as gains vs. losses. & Re-articulate all options in neutral, dual-frame economic balances. \\ \hline
\textbf{Sunk Cost Fallacy} & Continuing failing projects because heavy funds or effort were already spent. & Zero-based budgeting: Evaluate project as if inheriting it today for the first time. \\ \hline
\textbf{Status Quo Bias} & Defaulting to existing operational routines due to loss aversion. & Invert default options; require affirmative justification for doing nothing. \\ \hline
\textbf{Overconfidence Bias} & Underestimating project completion times and error probabilities. & Perform a Project Pre-Mortem: ``Assume we failed 1 year from now; why?'' \\ \hline
\textbf{Halo Effect} & Assuming an executive's technical idea is brilliant because of their general prestige. & Anonymize ideation voting; judge concepts blind to author identity. \\ \hline
\textbf{Dunning-Kruger Effect} & Incompetent performers drastically overestimating their domain mastery. & Standardized peer code/design reviews and objective rubric benchmarking. \\ \hline
\textbf{Groupthink} & Suppressing dissent in pursuit of harmonious consensus. & Silent Brainwriting before speaking; round-robin critique requirements. \\ \hline
\textbf{In-Group Favoritism} & Favoring solutions from one's own department over external options. & Cross-functional review squads; external user-testing panels. \\ \hline
\textbf{Automation Bias} & Blindly trusting algorithms or AI outputs without validation. & Mandate Human-in-the-Loop decision verification for critical thresholds. \\ \hline
\end{tabularx}
\end{frameworkbox}

\section{Multi-Criteria Decision Analysis (MCDA) Mathematical Model}

\begin{metricbox}{MCDA Weighted Scoring Formulation}
Given $n$ alternative candidate solutions $A = \{A_1, A_2, \dots, A_n\}$ evaluated against $m$ decision criteria $C = \{C_1, C_2, \dots, C_m\}$:
\[
S_i = \sum_{j=1}^m w_j \cdot r_{ij} \quad \text{subject to} \quad \sum_{j=1}^m w_j = 1.0, \quad w_j \ge 0
\]
Where:
\begin{itemize}
    \item $w_j$ is the normalized importance weight of criterion $C_j$.
    \item $r_{ij}$ is the normalized rating score of alternative $A_i$ on criterion $C_j$ ($r_{ij} \in [0, 10]$).
    \item $S_i$ is the aggregate performance index of alternative $A_i$. The optimal candidate satisfies $A^* = \arg\max_{i} S_i$.
\end{itemize}
\end{metricbox}

\section{Decision Tree Analysis \& Expected Value (EV)}

\begin{metricbox}{Expected Monetary / Utility Value Formula}
\[
\operatorname{EV}(d_k) = \sum_{s=1}^S P(s) \cdot V(d_k, s) - C(d_k)
\]
Where $P(s)$ is the probability of environmental state $s$, $V(d_k, s)$ is the gross payoff of decision $d_k$ under state $s$, and $C(d_k)$ is the upfront execution cost.
\end{metricbox}

\section{The MECE Principle (Mutually Exclusive, Collectively Exhaustive)}

\begin{frameworkbox}{Barbara Minto's MECE Decomposition Rule}
\begin{itemize}
    \item \textbf{Mutually Exclusive (ME):} No two categories overlap. Every discrete element belongs to exactly one category. Eliminates double-counting and boundary confusion.
    \item \textbf{Collectively Exhaustive (CE):} The sum of all categories covers 100\% of the problem space. No blind spots or unclassified outliers remain.
\end{itemize}
\end{frameworkbox}

\chapter{Unit 4: Ideation and Creative Problem Solving Toolkits}
\label{chap:unit4-ideation-selection}

\section{The SCAMPER Creative Heuristic Matrix}

\begin{canvasbox}{Alex Osborn \& Bob Eberle's SCAMPER Framework}
\begin{tabularx}{\linewidth}{|l|X|X|} \hline
\textbf{Letter} & \textbf{Core Operative Trigger} & \textbf{Analytical Engineering Question} \\ \hline
\textbf{S - Substitute} & Replace materials, components, personnel, energy sources. & What material or software library can replace the current bottleneck? \\ \hline
\textbf{C - Combine} & Merge disparate functions, services, workflows, hardware. & Can we combine user authentication with payment checkout in 1 step? \\ \hline
\textbf{A - Adapt} & Import solutions from nature, biology, or other industries. & How does the airline industry solve standby overbooking? \\ \hline
\textbf{M - Modify / Magnify} & Exaggerate, scale up, shrink, change color, change cadence. & What happens if we magnify the interface 10x or reduce processing to 1s? \\ \hline
\textbf{P - Put to Other Use} & Repurpose residual outputs, waste energy, idle servers. & Can heat generated by servers warm office spaces? \\ \hline
\textbf{E - Eliminate} & Strip non-essential features, reduce complexity, remove forms. & What happens if we completely eliminate user passwords? \\ \hline
\textbf{R - Reverse / Rearrange} & Invert sequence, swap causes and effects, flip roles. & What if the customer grades the company before using the service? \\ \hline
\end{tabularx}
\end{canvasbox}

\section{Divergent Ideation Protocols}

\begin{frameworkbox}{1. Brainwriting 6-3-5 Protocol}
\begin{itemize}
    \item \textbf{Structure:} 6 participants, write 3 ideas each, passed around every 5 minutes for 6 rounds.
    \item \textbf{Yield:} Exactly $6 \times 3 \times 6 = 108$ distinct ideation variations generated in 30 minutes.
    \item \textbf{Pedagogical Advantage:} Eliminates loud-voice domination; guarantees equal cognitive contribution.
\end{itemize}
\end{frameworkbox}

\begin{frameworkbox}{2. Crazy Eights Divergence Sprint}
\textit{Fast sketching exercise where each participant folds a sheet of paper into 8 rectangles and sketches 8 distinct ideas in 8 minutes (1 minute per cell). Prevents self-censorship and forces breakthrough thinking.}
\end{frameworkbox}

\section{Solution Evaluation: Desirability, Feasibility, Viability (DFV)}

\begin{metricbox}{IDEO's Three Lenses of Innovation}
A viable solution exists at the central nexus where all three conditions are satisfied:
\begin{enumerate}[label=\textbf{\arabic*.}, leftmargin=2em]
    \item \textbf{Desirability (Human Lens):} Do users actually want, value, and emotionally embrace this solution?
    \item \textbf{Feasibility (Technical Lens):} Can it be engineered using accessible technology, architecture, and regulatory parameters?
    \item \textbf{Viability (Economic Lens):} Is the unit economics sustainable? Does it generate positive ROI or operational surplus?
\end{enumerate}
\end{metricbox}

\section{The Pugh Decision Matrix (Controlled Convergence)}

\begin{metricbox}{Stuart Pugh's Concept Selection Formula}
Select a known baseline (Concept $B$). Evaluate candidates across $k$ criteria:
\[
\text{Net Score} = \sum (+) - \sum (-)
\]
Where:
\begin{itemize}
    \item $(+)$ indicates candidate is superior to baseline on this criterion.
    \item $(-)$ indicates candidate is inferior to baseline on this criterion.
    \item $(0)$ indicates candidate performs equivalently to baseline.
\end{itemize}
Candidates with high negative net scores are eliminated; positive candidates are synthesized.
\end{metricbox}

\section{Prioritization \& Selection Matrices}

\begin{canvasbox}{The 2$\times$2 Impact vs. Effort Prioritization Matrix}
\begin{tabularx}{\linewidth}{|l|X|X|} \hline
 & \textbf{Low Effort (Fast / Cheap)} & \textbf{High Effort (Complex / Costly)} \\ \hline
\textbf{High Impact} & \textbf{Quick Wins:} Execute immediately. High payoff, low barrier. Builds team momentum. & \textbf{Strategic Bets / Major Projects:} Plan carefully, allocate senior engineers, decompose into iterative sprints. \\ \hline
\textbf{Low Impact} & \textbf{Fill-ins / Low-Hanging Fruit:} Execute during low-load intervals or assign as onboarding tasks. & \textbf{Money Pits / Thankless Tasks:} Avoid, deprioritize, or ruthlessly eliminate from roadmap. \\ \hline
\end{tabularx}
\end{canvasbox}

\begin{frameworkbox}{The MoSCoW Prioritization Rule}
\begin{itemize}
    \item \textbf{Must Have (M):} Non-negotiable regulatory, safety, or core functional requirements without which product cannot launch.
    \item \textbf{Should Have (S):} High-value features essential for competitive success, but temporary manual workarounds exist.
    \item \textbf{Could Have (C):} Desirable enhancements implemented only if surplus sprint capacity remains.
    \item \textbf{Won't Have (W):} Explicitly deferred out-of-scope features for current development horizon.
\end{itemize}
\end{frameworkbox}

\chapter{Unit 5: Prototyping, Workflow Design, and Validation}
\label{chap:unit5-prototyping-validation}

\section{The Multi-Tier Service Blueprint Architecture}

\begin{canvasbox}{G. Lynn Shostack's Service Blueprint Matrix}
\begin{tabularx}{\linewidth}{|l|X|} \hline
\textbf{Tier} & \textbf{Operational Components} \\ \hline
\textbf{Physical Evidence} & Receipts, signage, mobile screen UI, hardware enclosures, automated emails. \\ \hline
\textbf{Customer Actions} & Every discrete action performed by the user across the service timeline. \\ \hline
\multicolumn{2}{|c|}{\textbf{--- Line of Interaction ---}} \\ \hline
\textbf{Frontstage Contact} & Direct user-facing actions performed by service staff or interactive web kiosks. \\ \hline
\multicolumn{2}{|c|}{\textbf{--- Line of Visibility ---}} \\ \hline
\textbf{Backstage Contact} & Behind-the-scenes actions triggered by frontstage events (e.g., barista brewing coffee). \\ \hline
\multicolumn{2}{|c|}{\textbf{--- Line of Internal Interaction ---}} \\ \hline
\textbf{Support Processes} & Enterprise databases, third-party payment gateways, cloud microservices, logistics. \\ \hline
\textbf{Failure Points [!]} & High-risk failure nodes requiring automated fallback redundancies. \\ \hline
\textbf{Wait Points [T]} & Customer bottleneck zones requiring progress disclosure or queue mitigation. \\ \hline
\end{tabularx}
\end{canvasbox}

\section{Validation Experiment Typology}

\begin{frameworkbox}{Spectrum of MVP Archetypes}
\begin{tabularx}{\linewidth}{|l|X|l|} \hline
\textbf{Experiment Type} & \textbf{Mechanism} & \textbf{Validates} \\ \hline
\textbf{Smoke Test / Fake Door} & A landing page with product description and a ``Sign Up'' button before backend exists. & Market demand and customer intent. \\ \hline
\textbf{Wizard of Oz} & Front-end looks fully autonomous; back-end tasks are executed manually by humans. & User workflow, layout, and perceived value. \\ \hline
\textbf{Concierge MVP} & The service is delivered 100\% manually, explicitly disclosing human labor. & Deep qualitative needs and willingness to pay. \\ \hline
\textbf{Prototype Spike} & Bare minimum functional software built solely to test technical feasibility. & API latency, algorithm convergence, memory load. \\ \hline
\end{tabularx}
\end{frameworkbox}

\section{Quantitative Usability Metrics}

\begin{metricbox}{1. System Usability Scale (SUS) Score (Brooke, 1986)}
10-item Likert scale (1 = Strongly Disagree to 5 = Strongly Agree). Let $R_i$ be response to question $i$:
\[
\text{SUS Score} = 2.5 \times \left[ \sum_{i \in \text{odd}} (R_i - 1) + \sum_{i \in \text{even}} (5 - R_i) \right]
\]
Interpretation: $\text{SUS} \ge 68.0$ represents average industry benchmark; $\text{SUS} \ge 80.3$ indicates grade A usability.
\end{metricbox}

\begin{metricbox}{2. Task Completion Rate (TCR) \& Single Ease Question (SEQ)}
\[
\text{TCR} = \frac{\text{Number of successfully completed test tasks}}{\text{Total task attempts}} \times 100\%
\]
\[
\text{SEQ Score} = \text{Average single 1-to-7 rating on: ``Overall, how easy was this task?''}
\]
\end{metricbox}

\section{The Feedback Capture Grid}

\begin{canvasbox}{4-Quadrant Feedback Synthesis}
\begin{tabularx}{\linewidth}{|X|X|} \hline
\textbf{Likes / What Worked ($\oplus$):} & \textbf{Criticisms / Friction Points ($\triangle$):} \\
$\bullet$ Features that generated unprompted praise. & $\bullet$ Confusion, hesitation, error loops. \\
$\bullet$ Interactions perceived as fast or intuitive. & $\bullet$ Missing functionality or unclear affordances. \\ \hline
\textbf{Questions Asked (?):} & \textbf{New Ideas / Suggestions ($\star$):} \\
$\bullet$ Doubts regarding pricing, security, storage. & $\bullet$ Feature variations proposed by users. \\
$\bullet$ Ambiguities in labels or instructions. & $\bullet$ Unexpected use cases discovered during test. \\ \hline
\end{tabularx}
\end{canvasbox}

\section{Continuous Iteration Engines: PDCA \& Agile Sprints}

\begin{frameworkbox}{Deming's PDCA Continuous Improvement Cycle}
\[
\textbf{Plan} \xrightarrow{\text{Hypothesis}} \textbf{Do} \xrightarrow{\text{Prototype Pilot}} \textbf{Check} \xrightarrow{\text{Measure Telemetry}} \textbf{Act} \xrightarrow{\text{Standardize/Pivot}} \textbf{Plan}
\]
\end{frameworkbox}

\chapter{Unit 6: Solution Communication and Professional Reflection}
\label{chap:unit6-communication-reflection}

\section{Barbara Minto's Pyramid Principle}

\begin{frameworkbox}{Top-Down Executive Communication Architecture}
\begin{itemize}
    \item \textbf{Core Premise:} State the governing conclusion / recommendation first; group supporting arguments underneath hierarchically.
    \item \textbf{SCQA Introduction Framework:}
    \begin{itemize}
        \item \textbf{Situation (S):} Establish stable context and shared facts that all stakeholders readily accept.
        \item \textbf{Complication (C):} Introduce the problem, unexpected shift, friction, or crisis disrupting the situation.
        \item \textbf{Question (Q):} Explicitly state the central question arising from the complication: ``What should we do?''
        \item \textbf{Answer (A):} Deliver your proposed solution and core thesis immediately.
    \end{itemize}
\end{itemize}
\end{frameworkbox}

\section{The Executive Pitch Narrative Arc}

\begin{canvasbox}{7-Beat Persuasive Design Presentation Structure}
\begin{tabularx}{\linewidth}{|l|l|X|} \hline
\textbf{Beat} & \textbf{Presentation Segment} & \textbf{Core Objective} \\ \hline
\textbf{1} & \textbf{The Hook / Human Anchor} & Open with a visceral human story, startling user metric, or unaddressed friction point. \\ \hline
\textbf{2} & \textbf{Problem Framing \& Market Gap} & Contrast the broken status quo with user frustration; define the root cause. \\ \hline
\textbf{3} & \textbf{The Breakthrough Solution} & Unveil the product/service concept with demo, clear affordances, and unique value proposition. \\ \hline
\textbf{4} & \textbf{Validation Evidence \& Telemetry} & Showcase test data: SUS scores, task success rates, customer quotes, pilot conversion. \\ \hline
\textbf{5} & \textbf{Feasibility \& Architecture} & Outline technology stack, system integration, regulatory compliance, and risk mitigation. \\ \hline
\textbf{6} & \textbf{Unit Economics \& Roadmap} & Define budget requirements, milestone timeline, ROI, and resource allocation. \\ \hline
\textbf{7} & \textbf{The Call to Action (The Ask)} & State the unambiguous next decision or budget approval required from leadership. \\ \hline
\end{tabularx}
\end{canvasbox}

\section{Experiential Learning \& Reflective Practice Cycles}

\begin{frameworkbox}{1. David Kolb's Experiential Learning Cycle (1984)}
\begin{center}
$\text{Concrete Experience} \to \text{Reflective Observation}$ \\
$\downarrow$ \hfill $\uparrow$ \\
$\text{Active Experimentation} \leftarrow \text{Abstract Conceptualization}$
\end{center}
\begin{itemize}
    \item \textit{Concrete Experience:} Doing the engineering task, encountering real design failure.
    \item \textit{Reflective Observation:} Stepping back to examine what happened and what discrepancies emerged.
    \item \textit{Abstract Conceptualization:} Formulating general principles, new mental models, and hypotheses.
    \item \textit{Active Experimentation:} Testing the new mental model in the next design iteration.
\end{itemize}
\end{frameworkbox}

\begin{canvasbox}{2. Graham Gibbs' Reflective Cycle (1988)}
\begin{enumerate}[label=\textbf{\arabic*.}, leftmargin=2em]
    \item \textbf{Description:} What happened? (Objective chronological narrative without judgment).
    \item \textbf{Feelings:} What were you thinking and feeling during the crisis/milestone?
    \item \textbf{Evaluation:} What was good and bad about the experience?
    \item \textbf{Analysis:} What sense can you make of the situation? (Apply design theories and root causes).
    \item \textbf{Conclusion:} What else could you have done? What did you personally learn?
    \item \textbf{Action Plan:} If it arose again, what specific protocol would you execute differently?
\end{enumerate}
\end{canvasbox}

\section{Blameless Postmortem Architecture}

\begin{frameworkbox}{Engineering Retrospective Framework (John Allspaw Heuristics)}
\begin{itemize}
    \item \textbf{Core Axiom:} Human error is the \textit{symptom} of a deeper system flaw, not the cause.
    \item \textbf{Timeline Reconstruction:} Second-by-second chronicle of system telemetry, human alerts, and actions taken.
    \item \textbf{Contributing Systemic Factors:} Gaps in documentation, test coverage, alerting thresholds, fatigue.
    \item \textbf{Corrective Action Items:} Assign owner, priority (P0/P1), ticket number, and hard completion date. Never prescribe ``be more careful.''
\end{itemize}
\end{frameworkbox}


\end{document}
